\documentclass[10pt]{article}
\usepackage[T1]{fontenc}
\usepackage[utf8]{inputenc}
\usepackage[french]{babel}
\usepackage{amsmath,amssymb,mathrsfs,mathpazo}
\usepackage{geometry}
\geometry{a4paper,margin=2cm}
\usepackage{graphicx,array,tabularx,tabularray,colortbl,multirow,multicol,booktabs,enumitem,pifont,ulem,fancybox,diagbox,cancel}
\usepackage{pstricks,pstricks-add,pst-plot,pst-tree,pst-node,pst-text,pst-func,pst-math,pst-bspline,pst-3dplot,pst-eucl,pst-solides3d}
\let\pscalculate\undefined
\usepackage{pst-bezier}
\usepackage{tikz,pgfplots,tkz-tab}
\usetikzlibrary{arrows,arrows.meta,patterns,calc,positioning,trees,decorations.markings,shapes.geometric,shapes.misc,shapes.arrows,intersections,decorations.pathmorphing,decorations.pathreplacing,backgrounds,shadows,babel}
\pgfplotsset{compat=1.18}
\usepgfplotslibrary{fillbetween}
\usepackage[np]{numprint}
\usepackage{listings,makecell,pas-tableur,twoopt}
\let\barre\undefined
\usepackage[pstricks]{tablvar}
\definecolor{darkgreen}{RGB}{0,100,0}
\providecommand{\red}{\color{red}}
\providecommand{\blue}{\color{blue}}
\providecommand{\green}{\color{green}}
\providecommand{\black}{\color{black}}
\providecommand{\white}{\color{white}}
\providecommand{\gray}{\color{gray}}
\providecommand{\cyan}{\color{cyan}}
\providecommand{\magenta}{\color{magenta}}
\providecommand{\yellow}{\color{yellow}}
\providecommand{\violet}{\color{violet}}
\providecommand{\orange}{\color{orange}}
\providecommand{\euro}{\texteuro}
\providecommand{\EUR}{\texteuro}
\providecommand{\ieme}{\textsuperscript{e}}
\providecommand{\degres}{\ensuremath{^\circ}}
\providecommand{\degree}{\ensuremath{^\circ}}
\providecommand{\vv}[1]{\overrightarrow{#1}}
\providecommand{\Coord}[1]{#1}
\providecommand{\decosix}{}
\providecommand{\starredbullet}{\ensuremath{\bullet}}
\pagestyle{empty}
\setlength{\parindent}{0pt}
\usepackage{adjustbox}\DeclareUnicodeCharacter{2236}{:}\DeclareUnicodeCharacter{2190}{\ensuremath{\leftarrow}}
\providecommand{\R}{\mathbb{R}}
\providecommand{\N}{\mathbb{N}}
\providecommand{\Z}{\mathbb{Z}}
\providecommand{\Q}{\mathbb{Q}}
\providecommand{\D}{\mathbb{D}}
\providecommand{\C}{\mathbb{C}}
\providecommand{\e}{{\mathrm{e}}}
\providecommand{\vect}[1]{\overrightarrow{#1}}
\providecommand{\vectt}[1]{\overrightarrow{\mathrm{#1}}}
\providecommand{\vvt}[1]{\overrightarrow{\mathrm{#1}}}
\providecommand{\np}[2][]{{#2~#1}}
\providecommand{\vv}[1]{{\overrightarrow{#1}}}
\renewcommand{\d}{{\mathrm{d}}}
\providecommand{\degres}{{^\circ}}
\providecommand{\degre}{{^\circ}}
\providecommand{\blue}{\color{blue}}
\providecommand{\red}{\color{red}}
\providecommand{\enskip}{\quad}
\providecommand{\llbracket}{\lbrack\!\lbrack}
\providecommand{\rrbracket}{\rbrack\!\rbrack}
\providecommand{\euro}{{\text{€}}}
\providecommand{\no}{{\text{n°}}}
\providecommand{\interval}[2]{{[#1;#2]}}
\providecommand{\intervFO}[2]{{[#1;#2[}}
\providecommand{\dotfill}{\ldots}

\providecommand{\ds}{\displaystyle}
\providecommand{\textding}[1]{#1}
\providecommand{\Oij}{\(\mathrm{O};\vec{\imath},\vec{\jmath}\)}
\providecommand{\Oijk}{\(\mathrm{O};\vec{\imath},\vec{\jmath},\vec{k}\)}
\providecommand{\Ouv}{\(\mathrm{O};\vec{u},\vec{v}\)}
\providecommand{\intervFF}[2]{\ensuremath{[#1;#2]}}
\renewcommand{\euro}{\ifmmode\text{\texteuro}\else\texteuro\fi}
\DeclareUnicodeCharacter{00B0}{\ifmmode^{\circ}\else\textdegree\fi}
\DeclareUnicodeCharacter{20AC}{\ifmmode\text{\texteuro}\else\texteuro\fi}
% Énoncés : Éducation nationale. Transcription et corrigés : APMEP, Denis Vergès et contributeurs. Découpage et classement : CoopMaths, Sébastien Lozano, Éric Elter et contributeurs. Adaptation de la présentation : Maths Explorer.
% CC BY-SA : https://coopmaths.fr/www/about/
% https://coopmaths.fr/alea/static/e3c/2020/tex/e3c_2020_00_sujet4_3_cor.tex
\begin{document}
\noindent\textit{APMEP / CoopMaths. CC BY-SA. Présentation adaptée par Maths Explorer.}\par\medskip
\section*{Exercice 3 (5 points)}
	
	\subsection*{Partie A}
	
	\subsubsection*{1. Calculer $P(B \cap S)$}
	
	On peut dresser un arbre pondéré de probabilités :
	
	\begin{tikzpicture}
	[level 1/.style={level distance=2cm,
		sibling distance=3cm},
	level 2/.style={level distance=2.5cm,
		sibling distance=1.5cm}]
	\node {} [grow'=right]
	child {node {$S$}
		child {node {$B$}
			edge from parent node[above] {$\dfrac34$}
		}
		child {node {$\overline B$}
			edge from parent node[below] {$\dfrac14$}
		}
		edge from parent node[above] {$\dfrac12$}
	}
	child {node {$\overline S$}
		child {node {$B$}
			edge from parent node[above] {$\dfrac18$}
		}
		child {node {$\overline B$}
			edge from parent node[below] {$\dfrac78$}
		}
		edge from parent node[below] {$\dfrac12$}
	}
	;
\end{tikzpicture}

	
	On a $P(B \cap S) = P(S \cap B) = P(S) \times P_S(B) = \dfrac{1}{2} \times \dfrac{3}{4} = \dfrac{3}{8}$.
	
	\subsubsection*{2. Déterminer la probabilité que le groupe de copains réponde correctement à la question posée}
	
	D’après la loi des probabilités totales on a :
	\[
	P(B) = P(B \cap S) + P(B \cap \overline{S}) = \dfrac{1}{2} + \dfrac{1}{2} \times \dfrac{1}{8} = \dfrac{3}{8} + \dfrac{1}{16} = \dfrac{6}{16} + \dfrac{1}{16} = \dfrac{7}{16}
	\]
	
	\subsubsection*{3. Les évènements $S$ et $B$ sont-ils indépendants ?}
	
	On a avec $P(S) = \dfrac{1}{2}$ :
	\[
	P(B \cap S) = \dfrac{3}{8} \text{ et } P(B) \times P(S) = \dfrac{7}{16} \times \dfrac{1}{2} = \dfrac{7}{32}
	\]
	Les évènements $S$ et $B$ ne sont donc pas indépendants.
	
	\subsection*{Partie B}
	

	
	\subsubsection*{1. Déterminer la loi de probabilité de $X$}
	
	On a $P(\overline{B}) = 1 - \dfrac{7}{16} = \dfrac{9}{16}$. D’où le tableau :
	
	\[
	\begin{array}{|c|c|c|c|}
		\hline
		&&&\\
		X & -5 & 5 & 25 \\
		&&&\\
		\hline
		&&&\\
		P(X = x_i) & \dfrac{9}{16} & \dfrac{6}{16} & \dfrac{1}{16} \\
		&&&\\
		\hline
	\end{array}
	\]
	
	\subsubsection*{2. Que retourne la fonction \texttt{Jeu} écrite ci-dessous en langage Python avec les listes $L = [-5, 5, 25]$ et $G = [0,5625; 0,375; 0,0625]$ ?}
	
\begin{center}
	
	\begin{tabularx}{0.4\linewidth}{|X|}\hline
		def \texttt{Jeu(L,G):}\\
		\quad  \texttt{n = len(L)}\\
		\quad \texttt{E = 0}\\
		\quad \texttt{for i in range(n):}\\
		\qquad \texttt{E =E + L[i]*G[i]} \\
		\quad \texttt{return(E)}\\ \hline
	\end{tabularx}
\end{center}
	
	L’algorithme donne l’espérance mathématique du jeu :
	\[
	E = -5 \times \dfrac{9}{16} + 5 \times \dfrac{6}{16} + 25 \times \dfrac{1}{16} = \dfrac{-45 + 30 + 25}{16} = \dfrac{10}{16} = \dfrac{5}{8} = 0,625~\euro{}
	\]
	En moyenne ce jeu permettra de gagner 6,25€ en 10 parties.

\end{document}
